\documentclass[11pt]{article}

\usepackage[a4paper,margin=1in]{geometry}
\usepackage{amsmath,amssymb,amsthm,mathtools}
\usepackage{array,tabularx}

\newcommand{\Db}{\mathrm{D}^{-}(R)}
\newcommand{\XX}{\mathbf{X}}
\newcommand{\TT}{\mathcal{T}}
\newcommand{\thicktensor}[1]{\operatorname{thick}^{\otimes}\{#1\}}

\title{Proof Effort Summary}
\author{ShuYu}
\date{}

\begin{document}

\maketitle

\section{Problem Overview}

\noindent\textbf{Problem (Remark 7.22, Matsui--Takahashi 2017).}
Let $R$ be a commutative noetherian ring and $x \in R$ a nonzerodivisor. For a positive integer $d$ and non-negative integers $a_0,\ldots,a_d$, define
\[
X(a_0,a_1,\ldots,a_d)
=
\bigoplus_{i \geqslant 0}\left(R/x^{a_0 i^d + a_1 i^{d-1} + \cdots + a_d}\right)[i].
\]
Does the equality
\[
\thicktensor{X(a_0,\ldots,a_d)\mid a_0,\ldots,a_d \geqslant 0}
=
\thicktensor{X(1,0,\ldots,0)}
\]
hold for all $d \geqslant 3$? The case $d=2$ was proven in Example 7.21 of the same paper.

\paragraph{Classification.}
\begin{itemize}
\item \textbf{Area:} Commutative algebra / homological algebra; thick tensor ideals in $\Db$ and tensor triangular geometry.
\item \textbf{Type:} Equality of thick tensor ideals. The non-trivial direction is showing every generator on the left belongs to the single-generator ideal on the right.
\item \textbf{Difficulty:} Hard; an explicitly stated open research problem, unresolved for about nine years. The $d=2$ proof exploits a specific algebraic coincidence, namely $2$-fold even/odd splitting via Corollary 7.3, which does not directly generalize because $d!$ has odd prime factors for $d \geqslant 3$.
\end{itemize}

\section{Final Proof Status}

\noindent\textbf{Result: PASS. A correct proof was found.}

The proof establishes the affirmative answer for all $d \geqslant 1$, fully resolving Remark 7.22. The key strategy is:
\begin{enumerate}
\item \textbf{Generalize $2$-fold splitting to $p$-fold splitting (Lemma B):} for any $p \geqslant 2$, a complex $\XX(a_i)$ belongs to the thick tensor ideal generated by its $p$-fold sub-sampled parts $\XX(a_{pi+r})$, for $r=0,\ldots,p-1$. This overcomes the fundamental obstruction of dividing by odd primes.
\item \textbf{Bottom-up induction on monomial degree:} starting from $\XX(i^d)\in\TT$, iterate finite differences (Lemma D) to produce $\XX(\Delta^m(i^d))\in\TT$, subtract lower-order terms (Lemma C) using the inductive hypothesis, and apply $c$-fold splitting with $c=d!/k!$ to recover $\XX(i^k)$.
\end{enumerate}

The proof was verified with $78$ out of $79$ miniclaims passing, zero failures, and one UNCERTAIN point: a minor expository gap regarding the trivial case $k=d$ in the $c$-fold splitting step. This is mathematically inconsequential since $\XX(i^d)$ is already the generator of $\TT$.

\section{Round-by-Round Summary}

\subsection{Round 1 (Exploration)}

All three proof search agents, Claude, Codex, and Gemini, produced empty proofs with only placeholder text. However, the verification phase was still useful: Claude's verification report identified the fundamental algebraic obstruction. For $d \geqslant 3$, the finite difference cascade produces leading coefficients divisible by $d!$, which contains odd prime factors, while Corollary 7.3 of the paper only provides $2$-fold even/odd splitting, that is, division by $2$. This diagnosis shaped the subsequent strategy.

\paragraph{Verdict.}
FAIL. No proof content was produced.

\paragraph{Key output.}
Identification of the odd-prime obstruction.

\subsection{Round 2 (Breakthrough)}

Claude produced a complete, mathematically rigorous proof generalizing the result to all $d \geqslant 1$. The proof introduced four new lemmas, labeled A--D, and used bottom-up induction on monomial degree. The central innovation was Lemma B, the $p$-fold splitting lemma, which generalizes Corollary 7.3 from $p=2$ to arbitrary $p \geqslant 2$. Verification confirmed $22$ miniclaims, all PASS, before timing out. Codex and Gemini again produced empty proofs.

\paragraph{Verdict.}
FAIL, because verification was incomplete due to timeout.

\paragraph{Key output.}
The correct proof strategy was found.

\subsection{Round 3 (Refinement)}

All three models adopted Claude's proof strategy from Round 2. Gemini's version was selected because it had partial verification data, with $5$ out of $30$ miniclaims verified and all passing. The proofs added clarifications about $Y_r \in \Db$ and the $K$-flatness of free modules. Verification again timed out before the core mathematical content could be fully checked.

\paragraph{Verdict.}
FAIL, because verification was incomplete.

\paragraph{Key output.}
The proof exposition improved, but verification remained the bottleneck.

\subsection{Round 4 (Verification Success)}

The proof was further refined with explicit labeled properties $(\mathrm{T1})$--$(\mathrm{T5})$ of thick tensor ideals, detailed uniqueness proofs for the diagonal extraction in Lemma B, explicit well-definedness and injectivity checks in Lemma C, and a computational verification section. Both Claude and Codex received PASS verdicts with $78/79$ miniclaims passing. Claude was selected as the final proof.

\paragraph{Verdict.}
PASS.

\paragraph{Key output.}
Fully verified proof.

\section{Approaches Tried}

\begin{center}
\begin{tabularx}{\textwidth}{|>{\raggedright\arraybackslash}p{0.06\textwidth}|>{\raggedright\arraybackslash}p{0.34\textwidth}|>{\raggedright\arraybackslash}p{0.16\textwidth}|X|}
\hline
\textbf{\#} & \textbf{Strategy} & \textbf{Rounds} & \textbf{Outcome} \\
\hline
1 & No proof attempted; empty placeholders only & R1 (all models), R2 (Codex, Gemini) & Abandoned. The models failed to produce content in the early rounds. \\
\hline
2 & $p$-fold splitting + finite differences + bottom-up induction & R2--R4 (Claude, then all models) & Carried forward and verified. This was the winning strategy. \\
\hline
\end{tabularx}
\end{center}

Only one substantive proof strategy was attempted across all rounds. The key evolution was not in the mathematical approach but in the level of exposition: each round added more explicit justifications to help the verification agents complete their checks within the time budget.

\paragraph{Specific refinements across rounds.}
\begin{itemize}
\item \textbf{R2 to R3:} Added explicit justification for $Y_r \in \Db$ and $K$-flatness arguments for the derived tensor product.
\item \textbf{R3 to R4:} Added labeled properties $(\mathrm{T1})$--$(\mathrm{T5})$, a uniqueness proof for diagonal extraction, explicit injectivity verification, an iterated Lemma D proof, and computational verification for small cases.
\end{itemize}

\section{Key Mathematical Insights}

\subsection{Core Discovery: $p$-fold Splitting (Lemma B)}

The central insight that unlocks the general case is that Corollary 7.3's even/odd decomposition generalizes to arbitrary $p$-fold splitting. Given a complex
\[
X=\bigoplus_{n \geqslant 0} M_n[n]
\]
with zero differentials, define
\[
X^{(r)}=\bigoplus_{i \geqslant 0} M_{pi+r}[i]
\]
for $r=0,\ldots,p-1$. Then
\[
X \in \thicktensor{X^{(0)},\ldots,X^{(p-1)}}.
\]

The proof constructs an auxiliary complex
\[
Y_r=\bigoplus_{j \geqslant 0} R[(p-1)j+r],
\]
tensors it with $X^{(r)}$, and extracts the diagonal direct summand
\[
Z_r=\bigoplus_{i \geqslant 0} M_{pi+r}[pi+r].
\]
This is the tool that enables division by any integer $p$, overcoming the odd-prime obstruction.

\subsection{Coefficient Formula via Stirling Numbers (Lemma A)}

The $m$-th forward finite difference of $i^d$ has the explicit expansion
\[
\Delta^m(i^d)=\sum_{j=0}^{d-m}\binom{d}{j}\,m!\,S(d-j,m)\,i^j,
\]
where all coefficients are non-negative integers. This non-negativity is essential: it ensures that the lower-order remainder
\[
q(i)=\Delta^m(i^d)-\frac{d!}{k!}\,i^k
\]
has non-negative coefficients, enabling subtraction via Lemma C.

\subsection{Induction Architecture}

The proof proceeds by induction on $k$ from $0$ to $d$, showing $\XX(i^k)\in\TT$.
\begin{itemize}
\item \textbf{Base case ($k=0$):} Extract $R/x$ as a direct summand of $\XX(i^d)$, then tensor with $\bigoplus_{i \geqslant 0} R[i]$.
\item \textbf{Inductive step ($k \geqslant 1$):} Apply $\Delta^{d-k}$ to $i^d$ via Lemma D, subtract lower-order terms using Lemma C and the inductive hypothesis, and use $c$-fold splitting with $c=d!/k!$ via Lemma B to recover $\XX(i^k)$.
\end{itemize}

\subsection{Structural Observation}

The proof reveals that the case $d=2$ is not genuinely special. The same mechanism works uniformly for all $d$. The apparent obstruction at odd primes was an artifact of the limited toolkit, namely only having $2$-fold splitting, rather than a genuine mathematical barrier.

\end{document}
